\chapter{Syntax and Program Structure}\label{ch:syntax}

Dylan, as the 1992 manual defines it, is written in a parenthesized prefix
notation of the kind Lisp and Scheme use. A program is a sequence of
\emph{forms}; each form is either an atom, such as a number, a string or a
name, or a list whose first element says what the form does. There are no
statements, no operator precedence and no semicolons ending lines: a
semicolon starts a comment.

This chapter describes the notation itself: how the reader turns
characters into Dylan objects (Section~\ref{sec:syntax-reader}), how those objects
are evaluated as expressions (Section~\ref{sec:syntax-eval}), and how a program is
put together from top-level forms and source files
(Section~\ref{sec:syntax-program}). The standard is \dylanbook{chapters 2 and 3};
symbols and keywords are also covered in \dylanbook{chapter 15}. Both
implementations read programs with the same reader, so everything in
Section~\ref{sec:syntax-reader} holds for both.

\section{How the examples are shown}\label{sec:syntax-conventions}

The examples in this part of the book are Dylan source. Each has been run
with both implementations, compiled with \cmd{bin/dylanc} and typed at the
read-eval-print loop of the Common Lisp implementation, except where the
text says that an example applies to only one of them. The value of an
expression is shown in a comment after it, introduced by \code{=>}:

%! skip
\begin{dylan}
(+ 3 4)                 ; => 7
\end{dylan}

\noindent When an expression returns several values, the comment lists them
in order, separated by commas; \code{(values 98 102)} is shown as
\code{; => 98, 102}. When it returns no values at all, the comment says
\code{(no values)}. A long value continues on the comment lines below. Output that a form prints is shown in comments
introduced by \code{prints:}, one line of output to each comment:

%! skip
\begin{dylan}
(print "hello")         ; prints: "hello"
\end{dylan}

\noindent A form whose evaluation signals an error is marked
\code{; => error}, usually followed by a short description of the error.
The exact wording of error messages differs between the two
implementations, and the comment does not try to reproduce either.

Values that cannot be read back in, such as classes and functions, are
printed inside braces: \code{\{the class <integer>\}}, as the 1992
manual writes them.

\section{The reader}\label{sec:syntax-reader}

The \term{reader} turns the characters of a program into Dylan objects:
numbers, strings, symbols, lists and so on. A program is just those objects;
evaluation (Section~\ref{sec:syntax-eval}) gives them meaning. The reader is the
same in both implementations: the compiler to C uses the Common Lisp
implementation's reader, so a file that one reads, the other reads
identically.

\subsection{White space and delimiters}

Spaces, tabs, newlines, returns and page breaks are \term{white space}. White
space separates tokens, and the amount of it does not matter, so a program
can be laid out freely. A token also ends at a parenthesis, a double quote,
a semicolon, a quote mark, a backquote or a comma. Everything else is
part of a token, which is why names such as \code{set!}, \code{even?},
\code{<integer>} and \code{*standard-output*} are single names.

\subsection{Comments}\label{sec:syntax-comments}

A semicolon\index{comment} starts a comment that runs to the end of the
line. The characters \code{\#|} start a block comment, which runs to the
matching \code{|\#}. Block comments may span lines and may be nested, so a
block comment can comment out code that already contains one:

\begin{dylan}
(+ 1 #| an inline comment |# 2)   ; => 3
#| (print "never run")
   #| a nested comment |#
   (print "still inside the outer comment") |#
\end{dylan}

\noindent A comment is white space as far as the reader is concerned. Note
that \code{\#;}, which comments out one form in some Lisps, does not exist
in Dylan.

\subsection{Numbers}\label{sec:syntax-numbers}

Numbers are written in the usual way. Chapter~\ref{ch:numbers} describes the
number classes and arithmetic; here is only how numbers are written.

\begin{itemize}
\item An \term{integer} is an optional sign followed by decimal digits:
  \code{0}, \code{42}, \code{-17}, \code{+57}. A trailing decimal point is
  allowed and still denotes an integer: \code{1.} is the integer \code{1}.
\item A \term{ratio} is two integers separated by a slash, with no spaces:
  \code{1/3}, \code{-6/4}. Ratios are always reduced, and a ratio that
  reduces to a whole number is an integer: \code{4/2} reads as \code{2}.
\item A \term{floating-point number} has a decimal point with digits after
  it, an exponent, or both: \code{1.5}, \code{-4.0}, \code{.5},
  \code{1e10}, \code{2.5d2}. An exponent is introduced by \code{e} or
  \code{d} for a double float, the default, or by \code{f} or \code{s} for
  a single float; \code{1.5s0} is a single float.
\item An integer or ratio in another base is written \code{\#x} (hexadecimal),
  \code{\#o} (octal) or \code{\#b} (binary) followed by the digits, which
  may have a sign: \code{\#x1f4e}, \code{\#xFF}, \code{\#o-17},
  \code{\#b101}. Letters may be in either case.
\end{itemize}

\begin{dylan}
(list 123 +57 -4.0 1.5 #x1f4e)     ; => (123 57 -4.0 1.5 8014)
(list 1. 6/4 4/2 .5)               ; => (1 3/2 2 0.5)
(list #o-17 #b101 #xff/2 1.0e-3)   ; => (-15 5 255/2 0.001)
\end{dylan}

\noindent Both implementations print double floats the same way, but they
print a single float and a large or small double float differently: the
compiler to C prints \code{1.5s0} and \code{10000000000.0} where the Common
Lisp implementation prints \code{1.5f0} and \code{1.0e10}.

\begin{warnbox}
The compiler to C has no bignums: an integer is 63 bits, and an integer
literal that does not fit is an error when compiling. It has no complex
numbers either. The Common Lisp implementation has both, and reads a
complex number in Common Lisp's notation, \code{\#c(1 2)}. See
Chapter~\ref{ch:numbers}.
\end{warnbox}

\subsection{Characters}\label{sec:syntax-characters}

A \term{character} is written \code{\#\textbackslash} followed by the
character itself, or by the character's name. The names of the 1992 manual
are \code{newline}, \code{space}, \code{rubout}, \code{page}, \code{tab},
\code{backspace}, \code{return} and \code{linefeed}; their case does not
matter, so \code{\#\textbackslash Space} and \code{\#\textbackslash space}
are the same character. \code{\#\textbackslash linefeed} is the same
character as \code{\#\textbackslash newline}. A single letter after the
backslash is that letter, in its case: \code{\#\textbackslash a} and
\code{\#\textbackslash A} are different characters. Any other single
character may follow the backslash too, parentheses, semicolons and quotes
included.

\begin{dylan}
(list #\a #\A #\( #\; #\")   ; => (#\a #\A #\( #\; #\")
(list #\Space #\TAB)         ; => (#\space #\tab)
(map (method (c) (as <integer> c))
     (list #\newline #\linefeed #\space #\tab #\rubout))
                             ; => (10 10 32 9 127)
\end{dylan}

\noindent Because the reader is built on Common Lisp's, it also accepts the
other character names Common Lisp knows, such as \code{\#\textbackslash
Nul} and \code{\#\textbackslash Escape}. A program that is to be portable to
other Dylan implementations should keep to the manual's names.

\subsection{Strings}\label{sec:syntax-strings}

A \term{string} is written between double quotes. Every character between
the quotes is part of the string, white space and newlines included. A
backslash quotes the character after it, which is how a string includes a
double quote or a backslash:

\begin{dylan}
"abc"              ; => "abc"
(size "a\"b")      ; => 3
(size "a\\b")      ; => 3
(element "a\"b" 1) ; => #\"
\end{dylan}

\noindent A backslash does nothing more than quote the next character, as
\dylanbook{chapter 3} says. In particular there are no escape sequences:
\code{"a\textbackslash nb"} is the three characters \code{a}, \code{n} and
\code{b}, not a string with a newline in it.

\begin{dylan}
(size "a\nb")      ; => 3
(element "a\nb" 1) ; => #\n
\end{dylan}

\noindent To put a newline into a string, write the newline itself between
the quotes, or build the string from characters at run time:

\begin{dylan}
(size (as <string> (list #\a #\newline #\b)))   ; => 3
\end{dylan}

\subsection{Symbols}\label{sec:syntax-symbols}

Any token that is not a number, and does not end in a colon, is a
\term{symbol}. Symbols name variables, and they are also objects in their
own right, of class \code{<symbol>}\index{<symbol>@\code{<symbol>}}. A
symbol may contain letters, digits and most punctuation: \code{+},
\code{<=}, \code{set!}, \code{even?}, \code{*x*}, \code{\$pi}, \code{1+},
\code{-} and \code{...} are all symbols.

\begin{dylan}
(list '+ '<= 'set! 'even? '*x* '$pi '1+ '- '...)
                       ; => (+ <= set! even? *x* $pi 1+ - ...)
(instance? 'foo <symbol>)   ; => #t
\end{dylan}

\noindent Symbols are not case-sensitive\index{case of symbols}:
\code{Foo}, \code{foo} and \code{FOO} are the same symbol. But a symbol
remembers the case it was first written in, and that is how it prints
(\dylanbook{chapter 15}). In the example below the symbol is written
\code{Hello} first, so it prints that way however it is written later:

\begin{dylan}
(id? 'Hello 'hello)        ; => #t
'HELLO                     ; => Hello
(as <string> 'hello)       ; => "Hello"
\end{dylan}

\noindent ``First written'' means first read. The compiler to C reads a
whole program before running any of it, so a symbol's spelling is the one
in the source text that comes first. The Common Lisp implementation reads a
form at a time, so a symbol first made at run time by
\code{(as <symbol> "Abc")} keeps the spelling \code{Abc} if no form read
earlier wrote it.

Dylan programs follow naming conventions (\dylanbook{chapter 3}). Except
for setter variables, described below, the conventions mean nothing to the
language; they only help a reader.

\begin{itemize}
\item Names of classes begin and end with angle brackets: \code{<window>},
  \code{<integer>}.
\item Module variables meant to change while the program runs begin and
  end with asterisks: \code{*window-count*}.
\item Constants begin with a dollar sign: \code{\$pi}.
\item Predicates, functions that answer true or false, end with a question
  mark: \code{even?}, \code{instance?}.
\item Operations that destructively modify their arguments end with an
  exclamation point: \code{reverse!}, \code{sort!}.
\end{itemize}

\subsection{Setter names}\label{sec:syntax-setter-names}

A variable name is either a symbol or a list of the form
\code{(setter \meta{symbol})}. The second kind names a \term{setter
variable}, the variable that holds the function which stores into the
place the function in \meta{symbol} reads. The value of \code{element} reads
an element of a collection; the value of \code{(setter element)} stores
one. A setter name can be used wherever a symbol can name a variable:

\begin{dylan}
(define (setter my-variable) 20)
(setter my-variable)            ; => 20
\end{dylan}

\noindent Because \code{setter} has this special meaning, the symbol
\code{setter} itself cannot be the name of a variable. The design note that
later replaced \code{(setter foo)} with \code{foo-setter} is not part of
this language. Setters are described in Section~\ref{sec:control-assignment} and
Section~\ref{sec:functions-setters}.

\subsection{Keywords}\label{sec:syntax-keywords}

A token ending in a colon is a \term{keyword}: \code{foo:}, \code{size:},
\code{init-value:}. Keywords look like symbols but are a separate class,
\code{<keyword>}\index{<keyword>@\code{<keyword>}}, and are used as
symbolic constants, chiefly to name keyword arguments. The colon is not
part of a keyword's name. Keywords, like symbols, are not case-sensitive and
remember the case they were first written in.

\begin{dylan}
foo:                         ; => foo:
(id? hello: HELLO:)          ; => #t
(instance? foo: <keyword>)   ; => #t
(instance? foo: <symbol>)    ; => #f
(as <keyword> "size")        ; => size:
(as <string> size:)          ; => "size"
\end{dylan}

\noindent Two keywords, or two symbols, with the same name are always the
same object: \code{id?} is true of them.

\subsection{True, false and the empty list}\label{sec:syntax-booleans}

\code{\#t}\index{t@\code{\#t}} is the canonical true value and
\code{\#f}\index{f@\code{\#f}} is the false value; either may be written in
upper case. \code{\#f} is the \emph{only} false value: every other object
counts as true when tested, including \code{0}, the empty string and the
empty list.

The empty list, written \code{()}, is an object of its own, of class
\code{<empty-list>}. Unlike in older Lisps it is not the same object as
\code{\#f}, and it is true:

\begin{dylan}
(id? '() #f)                    ; => #f
(if '() 'true 'false)           ; => true
(instance? '() <list>)          ; => #t
\end{dylan}

\subsection{Lists and vectors}\label{sec:syntax-lists}

A list is written as its elements between parentheses. A dot between the
last two elements writes a dotted list, whose final tail is not the empty
list. A vector, a one-dimensional array, is written like a list with a
\code{\#} before the opening parenthesis.

\begin{dylan}
'(a (b c) ())        ; => (a (b c) ())
'(a . b)             ; => (a . b)
'#(1 2 3)            ; => #(1 2 3)
(instance? '#(1 2) <vector>)   ; => #t
\end{dylan}

\noindent A list in a program is normally a form to be evaluated, so a list
meant as data must be quoted. A vector need not be: \code{\#(1 2 3)}
evaluates to itself. The elements of a literal list or vector are not
evaluated, and literal data should not be modified.

\subsection{Quote and backquote}\label{sec:syntax-quote}

A single quote before a form, \code{'\meta{form}}, is an abbreviation for
\code{(quote \meta{form})}\index{quote@\code{quote}}, which evaluates to
the form itself, unevaluated. Quoting is needed for symbols and lists used
as data; other literals evaluate to themselves and need no quote, though
quoting them does no harm.

\begin{dylan}
(quote foo)          ; => foo
'foo                 ; => foo
'(+ 10 10)           ; => (+ 10 10)
(+ 10 10)            ; => 20
'"s"                 ; => "s"
\end{dylan}

A \term{backquote} quotes a template in which parts are filled in when it is
evaluated. Within a backquoted form, a comma, \code{,\meta{form}},
evaluates \meta{form} and puts its value in that place; \code{,@\meta{form}}
evaluates \meta{form}, which must answer a list, and splices its elements
in. \dylanbook{chapter 3} defers backquote to a section the manual does not
contain, so both implementations follow the usual Lisp rules, including
backquotes nested inside backquotes and backquoted vectors.

\begin{dylan}
(define x 5)
(define xs '(a b))
`(x ,x)              ; => (x 5)
`(1 ,@xs 2)          ; => (1 a b 2)
`(1 . ,x)            ; => (1 . 5)
`#(1 ,x)             ; => #(1 5)
\end{dylan}

\noindent The reader turns \code{'x}, \code{`x}, \code{,x} and \code{,@x}
into the lists \code{(quote x)}, \code{(quasiquote x)}, \code{(unquote x)}
and \code{(unquote-splicing x)}. The compiler to C prints these lists in
that long form; the Common Lisp implementation prints them abbreviated,
as \code{'x} and so on. A comma outside any backquote is an error. Backquote
is most useful in macros, an extension of the compiler to C described in
Chapter~\ref{ch:extensions}.

\subsection{\#key, \#rest and \#next}\label{sec:syntax-parameter-tokens}

The tokens \code{\#key}\index{key@\code{\#key}},
\code{\#rest}\index{rest@\code{\#rest}} and
\code{\#next}\index{next@\code{\#next}} mark the parts of a parameter
list (Section~\ref{sec:functions-parameters}); \code{\#rest} also appears in
\code{bind} (Section~\ref{sec:control-bind}). Like \code{\#t} and \code{\#f} they are
case-insensitive.

Both implementations also read \code{!key}, \code{!rest} and \code{!next}
as these same tokens, because Thomas, from which the Common Lisp
implementation grew, wrote them that way, and so do its programs, the
manual's examples among them. As a consequence those three names cannot
name variables. Any other name may.

\subsection{Unreadable objects}

Some objects, such as functions and classes, have no notation that reads
back as the same object. They print inside braces, \code{\{the generic
function print\}}, and there is no notation that reads them back. Any other
\code{\#} syntax than those above (\code{\#t}, \code{\#f},
\code{\#\textbackslash}, \code{\#(}, \code{\#x}, \code{\#o}, \code{\#b},
\code{\#|}, \code{\#key}, \code{\#rest}, \code{\#next}, and \code{\#c} in the
Common Lisp implementation) is an error.

\section{Expressions and evaluation}\label{sec:syntax-eval}

A program is made of expressions. Evaluating an expression yields zero or
more values, may have side effects, and may evaluate other expressions
inside it. There are four kinds of expression (\dylanbook{chapter 3}):
literal constants, variable references, function calls and syntax forms.

\subsection{Literal constants}

Numbers, characters, strings, keywords, vectors, \code{\#t}, \code{\#f}
and the empty list \code{()} evaluate to themselves. Symbols and non-empty
lists must be quoted to be used as constants, because unquoted they are
variable references and calls.

\begin{dylan}
"abc"      ; => "abc"
123        ; => 123
foo:       ; => foo:
#\a        ; => #\a
#t         ; => #t
()         ; => ()
'(1 2 3)   ; => (1 2 3)
\end{dylan}

\subsection{Variable references}

A symbol, or a setter name \code{(setter \meta{symbol})}, evaluates to the
value of the variable it names. There are two kinds of variable.
\term{Lexical variables} are created by \code{bind}, by a method's parameter
list and by several other forms, and are visible only in the text of the
form that creates them. \term{Module variables} are created by the defining
forms (Section~\ref{sec:syntax-program}) and are visible throughout their module
(Chapter~\ref{ch:modules}). An inner variable \emph{shadows} any outer one of the
same name, lexical or module.

\begin{dylan}
(define my-variable 25)
my-variable                     ; => 25
(bind ((my-variable 50))
  (+ my-variable my-variable))  ; => 100
\end{dylan}

\noindent Referring to a variable that has no value is an error. The
compiler to C also warns, when compiling, of a module variable that is used
but defined nowhere in the program.

\subsection{Function calls}\label{sec:syntax-calls}

A list whose first element is not the name of a syntax operator is a
\term{function call}. Its first element, the \term{operator}, is evaluated
to give the function; the rest are evaluated to give the arguments, and the
function is called with them. The values of the call are the values the
function returns.

The operator may be any expression that evaluates to a function, not only a
name. In this Dylan is like Scheme and unlike Common Lisp:

\begin{dylan}
(+ 3 4)              ; => 7
(* (+ 3 4) 5)        ; => 35
((if #t + *) 4 5)    ; => 9
((method (x) (* x x)) 6)   ; => 36
\end{dylan}

\noindent The arguments are evaluated from left to right. The manual allows
the operator to be evaluated at any time before, during or after the
arguments, and the two implementations differ: the compiler to C evaluates
the operator first, and the Common Lisp implementation evaluates it last. A
program should not depend on either order.

%! pre: (define trail '())
%! pre: (define-method note (x) (set! trail (cons x trail)) x)
\begin{dylan}
(list (note 1) (note 2) (note 3))   ; => (1 2 3)
(reverse trail)                     ; => (1 2 3)
\end{dylan}

\noindent (Here \code{note} records each argument in a list as it is
evaluated.)

\subsection{Syntax forms}\label{sec:syntax-forms}

A list whose first element names a \term{syntax operator} is a
\term{syntax form}. A syntax form looks like a call, but it is not one:
each syntax operator has its own rule for which of its parts are evaluated,
when, and what the others mean. In \code{(bind ((x 10)) (+ x 1))}, for
example, \code{((x 10))} is not evaluated at all; \code{bind} reads it as a
list of bindings.

There are three kinds of syntax operator.

\begin{description}
\item[Special forms] are built into the language:
  \code{begin}, \code{bind}, \code{bind-exit}, \code{bind-methods},
  \code{if}, \code{method}, \code{quote}, \code{set!}, \code{setter} and
  \code{unwind-protect}.
\item[Macros] are syntax forms defined in terms of others. The language's
  macros are \code{and}, \code{or}, \code{when}, \code{unless}, \code{cond}, \code{case}, \code{select}, \code{for},
  \code{for-each}, \code{dotimes}, \code{while}, \code{until}, and the
  condition forms \code{handler-bind} and \code{handler-case}
  (Chapter~\ref{ch:conditions}). Backquote, \code{quasiquote}, is one too.
\item[Definitions] create module variables: \code{define},
  \code{define-class}, \code{define-generic-function}, \code{define-method},
  \code{define-slot} and \code{define-module}.
\end{description}

\noindent \code{not}, by contrast, is an ordinary function. A syntax operator is
not a function and has no value: \code{(map and ...)} is meaningless, and
so is passing \code{if} as an argument.

The 1992 manual has no way for a program to define its own syntax forms.
The compiler to C adds Common Lisp--style macros, defined with
\code{define-macro} and expanded when compiling (Chapter~\ref{ch:extensions}); the
Common Lisp implementation does not have them.

\subsection{Multiple values}

An expression can return one value, several, or none. \code{values} returns
its arguments as multiple values, and \code{bind} receives them
(Section~\ref{sec:control-values}). Where only one value is wanted, as an argument
to a call, for instance, the first value is used, or \code{\#f} if there are
none.

\begin{dylan}
(values 1 2 3)             ; => 1, 2, 3
(print (list (values 1 2) 3))   ; prints: (1 3)
(print (list (values)))         ; prints: (#f)
\end{dylan}

\section{Program structure}\label{sec:syntax-program}

\subsection{Top-level forms and definitions}

A form that is not part of any other form is a \term{top-level form}; so are
the forms directly inside a \code{begin} that is itself at top level.
A program is a sequence of top-level forms. Some are definitions, which
create module variables; the rest are expressions, evaluated for their
effects.

\begin{center}
\begin{tabular}{@{}lp{0.55\textwidth}@{}}
\toprule
Definition & Creates \\
\midrule
\code{define} & a variable with an initial value (Section~\ref{sec:control-define}) \\
\code{define-method} & a method, added to a generic function (Section~\ref{sec:functions-define-method}) \\
\code{define-generic-function} & a generic function (Section~\ref{sec:functions-define-generic}) \\
\code{define-class} & a class, its slots and their accessors (Section~\ref{sec:classes-define-class}) \\
\code{define-slot} & a slot of an existing class or object (Section~\ref{sec:classes-define-slot}) \\
\code{define-module} & a module (Chapter~\ref{ch:modules}) \\
\bottomrule
\end{tabular}
\end{center}

\noindent The compiler to C adds three more definitions, which are
described in Chapter~\ref{ch:extensions}: \code{define-macro},
\code{define-for-macros} and \code{define-thread-variable}.

The 1992 manual describes definitions as if they were ordinary forms that
could appear anywhere, and both implementations accept \code{define} inside
a method, where it still defines a module variable. Later design notes
(``Definitions are Declarative'') restrict definitions to top level, and
that is the only place they belong: a definition inside a method happens
again each time the method runs, and in the compiler to C a
\code{define-method} that refers to local variables cannot be compiled:
it stops with \code{define-method: the method refers to local variables
(k)}, naming the variables.
Put definitions at top level.

\subsection{When things happen}\label{sec:syntax-when}

The two implementations agree on what a program means, but not on when a
definition takes effect, because one runs a static program and the other
is interactive.

The \textbf{Common Lisp implementation} evaluates each top-level form when
it reaches it, at the prompt or in a file. A definition takes effect then,
and not before; a later form may redefine a variable, which only draws a
warning, or redefine a class or add methods to a generic function that is
already in use.

The \textbf{compiler to C} compiles a whole program at once. Every class,
generic function and method the program defines exists before any of its
expressions run, wherever in the program it is defined. The top-level
expressions, and the initial values of \code{define}, are then evaluated in
order, from the first file to the last. A method may therefore be called
by an expression that comes before its definition in the source, but a
variable defined with \code{define} has no value until its definition has
been evaluated.

%! c-only
\begin{dylan}
(later 2)                   ; => 20
(define-method later (x) (* x 10))
later-var                   ; => error: later-var has no value yet
(define later-var 5)
\end{dylan}

\noindent In the Common Lisp implementation the first expression is an error
too, since \code{later} is not yet defined when it is evaluated.

A static program is also stricter about definitions. In the compiler to C a
class or generic function may be defined only once, a name cannot be
defined both by \code{define} and as a generic function, and a variable
that holds a class or a generic function cannot be assigned. These are
errors when compiling. Chapter~\ref{ch:dylanc} lists the full set of restrictions.

\subsection{Source files and the module: header}\label{sec:syntax-files}

A Dylan source file conventionally has the extension \file{.dyl}. It may
begin with a \term{header} (design note 33): lines of the form
\meta{keyword}\code{:} \meta{value}, each starting at the beginning of a
line, ended by a blank line. The keyword \code{module:} names the module
the file's forms belong to; any other keywords, such as \code{author:},
\code{copyright:} and \code{version:}, are accepted and ignored. A file
without a header belongs to the module \code{dylan-user}.

%! skip
\begin{dylan}
module: shapes
author: A. Programmer
version: 1.0

(define-method greet ()
  (print "hello from shapes"))
(greet)
\end{dylan}

\noindent The header is not Dylan syntax: it is read before the reader
starts, and design note 33 does not allow comments in it. The module it names must be
defined by a \code{define-module} form, which can be in an earlier file of
the program; Chapter~\ref{ch:modules} describes modules. Both implementations
honor the header when they compile a file: the compiler to C for each file
on its command line, and the Common Lisp implementation in
\code{dylan:compile-file}. At the interactive prompt, \code{(dylan:in-module
\meta{name})} sets the module instead (Chapter~\ref{ch:lisp}).

A program for the compiler to C may be spread over several files, given on
the command line in order: \cmd{bin/dylanc -o app a.dyl b.dyl}. They form
one program, and their top-level expressions run in that order
(Chapter~\ref{ch:dylanc}).
